\section{Patristic substance and historical development}

The Fathers did not write a later Mariology manual in advance. Justin and
Irenaeus place Mary's obedient faith within the Eve--Mary recapitulation whose
saving head is Christ. Ignatius witnesses the reality and mystery of Christ's
birth from the Virgin. Athanasius and the Cappadocians protect the incarnate
Son's identity; Cyril's \emph{Theotokos} serves that Christology. Ephrem's
poetry praises holiness through a symbolic and liturgical register. Ambrose,
Jerome, Augustine, and later Greek homilists witness different parts of
virginity, holiness, ecclesial typology, and Marian destiny.

This evidence must be neither flattened nor dismissed. A Father can witness
the substance or matrix of a later dogma without using its final terms. A local
interpretation can also differ from the Church's mature judgment. Augustine's
famous exception for Mary concerns personal sin, not an explicit first-instant
Immaculate Conception. Thomas Aquinas taught sanctification after contraction
of original sin and therefore did not hold the 1854 formula, although his
Christology, account of grace, and universal mediation supply important
principles. Epiphanius's uncertainty about Mary's end prevents a claim of
patristic unanimity on her death.

Development is faithful when the same mystery is more exactly understood and
proposed; it is not proved by finding a translated word resembling a later
formula. Liturgy, controversy, episcopal consultation, theological argument,
and the sense of the faithful can witness and occasion development. None is an
independent revelation, and no opinion poll creates dogma.

